\subsection{Data Marts und Einsatz des KI-Modells zur KV-Kündiger Prognose}

Neben dem Training der KI-Modelle gibt es, wie zuvor beschrieben,
je Anwendungsfall auch eine Analysefragestellung, 
die den Einsatz des jeweiligen KI-Modells abbildet.
Dabei handelt es sich im Gegensatz zum Training
um strukturierte Datenanalysen.
Die Fragestellungen haben ein klar definiertes Ergebnis:
''Welche KV-Kunden und -Kundinnen sind gefährdet zu kündigen?'' 
oder ''Handelt es sich bei dieser Kfz-Schadensmeldung um einen Betrug?''.
Zusätzlich sind es sich wiederholende operative Analysen, 
die von den jeweiligen Sachbearbeitenden durchgeführt werden.
Dafür erhalten sie eine benutzerfreundliche UI, 
sodass keine hohe Datenkompetenz erforderlich ist.  

Eine Besonderheit bei diesen beiden Datenanalysen ist,
dass beim Einsatz der KI-Modelle offensichtlich 
die KI-Modelle selbst und damit indirekt auch ihre Trainingsdaten
genutzt werden.
Da es für die Modelle allerdings keine eigene Architekturkomponente
gibt, wird dieser Sachverhalt in Abb.~\ref{fig:datenarchitektur-regional}
jeweils angedeutet durch gerichtete Kanten 
von der Datenanalyse des Trainings eines KI-Modells
zur Datenanalyse, die seinen Einsatz darstellt.

Die aktuellen Daten, in denen das KI-Modell
kündigungsgefährdete KV-Kunden und Kundinnen identifizieren soll,
stammen alle aus strukturierten Datenquellen, 
ausgenommen sind lediglich die semi-strukturierten Daten
aus dem CRM-System.
Aber auch diese werden auf dem Weg vom Data Lake in das 
Data Warehouse in eine strukturierte Form gebracht.
Damit liegen alle benötigten Daten tagesaktuell im Data Warehouse vor.
Eine höhere Aktualität wird für die Analyse nicht benötigt.
Es ist ausreichend, wenn das Modell einmal täglich oder wöchentlich
eine Prognose über kündigungsgefährdete KV-Kunden und -Kundinnen macht.
Mit dieser Auswertung können die Mitarbeitenden dann diese Fälle 
besonders betreuen.
Häufiger als einmal täglich werden sich weder die 
Kündigungspotenziale ändern, noch können die Mitarbeitenden
häufiger darauf reagieren.
Demnach liegen alle benötigten Daten für die Analyse 
in ausreichender Aktualität im Data Warehouse vor.

Trotzdem werden die Daten nicht direkt aus dem Data Warehouse
gezogen, sondern es werden abhängige Data Marts dazwischengeschaltet.
Das hat den Grund, dass für die Analysefragestellung
nicht alle Daten aus dem Data Warehouse gebraucht werden.
Die zuständigen Sachbearbeitenden sollen nur die Informationen
erhalten, die für ihre Datenanalyse wichtig sind. 
Zusätzlich wird angenommen, dass die Betreuung der
privat Krankenversicherten durch die Mitarbeitenden
der zugehörigen Filiale passiert.
Deshalb werden in Abb.~\ref{fig:datenarchitektur-regional} 
viele Data Marts dargestellt.
Es soll je Filiale der Allianz in dieser Region einen eigenen 
Data Mart geben, der nur die Daten für diese Datenanalyse bereitstellt.
So erhalten die Mitarbeitenden einer Filiale auch nur 
die Kündigungsprognose für die Kunden und Kundinnen, 
die sie auch selbst betreuen. 
Und sie können keine personenbezogenen Daten
einsehen, die für sie irrelevant sind.

Von den Datenanalysen, die den Einsatz der KI-Modelle darstellen,
gibt es außerdem noch eine gerichtete Kante zurück zum Data Warehouse.
Diese symbolisiert, dass die Ergebnisse der Modelle für spätere Analysen 
abgespeichert werden.
Das ist beispielweise interessant, um die Performance der Modelle zu überwachen.